\documentclass[11pt,a4paper]{article}
\usepackage[utf8]{inputenc}
\usepackage[T1]{fontenc}
\usepackage{babel}
\renewcommand{\contentsname}{Table des mati\`eres}
\usepackage{lmodern}
\usepackage[margin=2cm]{geometry}
\usepackage{xcolor}
\usepackage{titlesec}
\usepackage{longtable}
\usepackage{booktabs}
\usepackage{colortbl}
\usepackage{graphicx}
\usepackage{parskip}
\usepackage{hyperref}
\usepackage{listings}
\usepackage{tcolorbox}
\tcbuselibrary{skins,breakable}
\usepackage{fancyvrb}

\definecolor{navy}{HTML}{111111}
\definecolor{gold}{HTML}{1565C0}
\definecolor{greytext}{HTML}{555555}
\definecolor{lightgrey}{HTML}{F2F4F9}
\definecolor{termbg}{HTML}{1E1E1E}
\definecolor{termgreen}{HTML}{6ECB63}
\definecolor{termgrey}{HTML}{D4D4D4}
\definecolor{codebg}{HTML}{F7F8FA}

\titleformat{\section}{\Large\bfseries\color{navy}}{\thesection}{0.6em}{}[{\color{navy}\titlerule}]
\titleformat{\subsection}{\large\bfseries\color{gold}}{\thesubsection}{0.6em}{}

\hypersetup{colorlinks=true, linkcolor=navy, urlcolor=navy, citecolor=navy}

\newcommand{\tblhead}[1]{\textcolor{white}{\textbf{#1}}}

\lstdefinestyle{code}{
  backgroundcolor=\color{codebg},
  basicstyle=\ttfamily\scriptsize,
  keywordstyle=\color{gold}\bfseries,
  commentstyle=\color{greytext}\itshape,
  stringstyle=\color{navy},
  showstringspaces=false,
  breaklines=true,
  frame=single,
  rulecolor=\color{lightgrey},
  framerule=0.6pt,
  xleftmargin=2pt,
  xrightmargin=2pt,
  aboveskip=6pt,
  belowskip=6pt,
  language=Python,
}
\lstset{style=code}

% --- Boite "capture terminal" : reconstitue fidelement (texte identique, copie-colle)
% la sortie reelle du terminal d'Ayman, presentee dans un cadre visuel type fenetre de
% terminal (barre de titre a puces) plutot qu'une image, puisqu'il ne s'agit pas d'une
% photo d'ecran mais du texte exact produit par les scripts, execute sur sa machine.
\newtcolorbox{terminal}[1][]{
  colback=termbg, colframe=termbg, colupper=termgrey,
  fonttitle=\bfseries\ttfamily\footnotesize,
  coltitle=white,
  boxrule=0pt, arc=3mm, top=8mm, bottom=2mm,
  breakable,
  overlay={
    \fill[red!70!black] (frame.north west) ++(4mm,-3.2mm) circle (1.1mm);
    \fill[yellow!80!orange] (frame.north west) ++(7.2mm,-3.2mm) circle (1.1mm);
    \fill[green!60!black] (frame.north west) ++(10.4mm,-3.2mm) circle (1.1mm);
  },
  #1
}
\fvset{fontsize=\footnotesize, formatcom=\color{termgrey}, xleftmargin=1mm}

\begin{document}

\begin{titlepage}
\centering
\vspace*{2cm}
{\Huge\bfseries\color{navy} Bilan du Sprint 2\par}
\vspace{0.5cm}
{\Large\bfseries\color{gold} Pipeline d'indexation complet -- Syst\`eme RAG pour hcp.ma\par}
\vspace{1.5cm}
{\large\color{greytext} Structure d'accueil~: Haut-Commissariat au Plan (HCP), Direction des Syst\`emes d'Information Statistiques\par}
\vspace{0.3cm}
{\large\color{greytext} Bin\^ome~: Ayman El Baida et Saad Mesbah\par}
\vspace{0.3cm}
{\large\color{greytext} P\'eriode couverte~: 14 -- 20 juillet 2026 (Sprint 2)\par}
\vspace{0.3cm}
{\large\itshape\color{greytext} Document r\'edig\'e le 20 juillet 2026 -- support de suivi pour l'encadrante et le rapport de stage\par}
\vfill
\end{titlepage}

\tableofcontents
\newpage

\section{Objectif du Sprint 2}
Le Sprint 1 avait \'etabli que le RAG s'appuie exclusivement sur les pi\`eces jointes PDF/XLSX/DOCX des publications hcp.ma (ADR 0003), avec un premier scraper capable de les t\'el\'echarger sur un \'echantillon fixe de 27 URLs choisies \`a la main. Le Sprint 2 (14 -- 20 juillet 2026) visait \`a transformer ces publications collect\'ees en donn\'ees r\'eellement interrogeables~: d\'ecoupage du texte en fragments (\emph{chunks}), calcul d'embeddings, indexation hybride s\'emantique~+~lexicale, et persistance en base SQLite -- ainsi qu'\`a lever une limite identifi\'ee en cours de sprint~: la liste fixe de 27 URLs ne pouvait pas suivre les nouvelles publications de hcp.ma dans le temps.

Cinq blocs de travail ont \'et\'e r\'ealis\'es~: (1) d\'ecouverte automatique des publications, (2) support des fichiers DOCX, (3) chunking et embeddings, (4) indexation hybride Chroma~+~BM25, (5) persistance SQLite et scripts d'orchestration. Chacun est d\'etaill\'e ci-dessous avec le code r\'eel, puis valid\'e par une ex\'ecution en conditions r\'eelles sur la machine d'Ayman (section~8).

\section{Vue d'ensemble du pipeline}
\begin{longtable}{@{}p{3.2cm}p{12.8cm}@{}}
\rowcolor{navy}
\tblhead{\'Etape} & \tblhead{R\^ole} \\
\midrule
\endhead
1. \texttt{Scraper} & Trouve les publications sur hcp.ma (d\'ecouverte automatique, ADR 0006) et t\'el\'echarge leurs pi\`eces jointes (PDF/XLSX/DOCX). \\
2. \texttt{Extracteur} & Extrait le texte brut et les tableaux de chaque pi\`ece jointe. \\
3. \texttt{IndexeurTexte} & D\'ecoupe le texte en chunks, calcule un embedding par chunk (BGE-M3), indexe dans Chroma (vectoriel) et BM25 (lexical). \\
4. \texttt{base\_donnees} & Ins\`ere document, chunks et indicateurs en base SQLite, avec d\'edoublonnage. \\
5. Recherche hybride & \`A partir d'une question, fusionne un classement s\'emantique (Chroma) et lexical (BM25) par \emph{Reciprocal Rank Fusion}. \\
\end{longtable}

\section{D\'ecouverte automatique des publications (ADR 0006)}

\subsection{Probl\`eme}
La liste fixe de 27 URLs du Sprint 1 ne peut ni couvrir l'historique complet des publications de hcp.ma, ni d\'etecter automatiquement une nouvelle publication. V\'erification faite sur le vrai site (\texttt{web\_fetch})~: chaque sous-th\`eme dispose d'une page ``listing'' pagin\'ee (ex. \texttt{Publications-Marche-du-travail\_r425.html}) qui liste ses publications, les plus r\'ecentes en premier, avec une pagination r\'eelle via \texttt{?start=N}.

\subsection{D\'ecision}
\texttt{Scraper.decouvrir\_urls\_liste} lit une page listing, en extrait les URLs d'articles r\'eels (filtr\'ees par motif \texttt{\_aNNNN.html} \emph{et} pr\'esence dans un titre \texttt{<h2>}--\texttt{<h5>}, pour \'eviter les faux positifs de menu), puis suit la pagination. Le param\`etre \texttt{max\_pages} distingue deux usages~: \texttt{max\_pages=1} (page la plus r\'ecente seulement, ex\'ecution r\'eguli\`ere de d\'etection de fra\^icheur) et \texttt{max\_pages=None} (toutes les pages, collecte historique compl\`ete).

\begin{lstlisting}[caption={\texttt{src/scraper.py} -- suivi de la pagination reelle}]
def decouvrir_urls_liste(self, url_listing, max_pages=None) -> list[str]:
    resp = requests.get(url_listing, headers=HEADERS, timeout=15)
    soup = BeautifulSoup(resp.text, "lxml")

    urls_trouvees = []
    for url in self._filtrer_urls_arabe(
        self._extraire_urls_articles(soup, url_listing)
    ):
        urls_trouvees.append(url)

    paliers = self._increments_pagination(soup)
    if max_pages is not None:
        paliers = paliers[: max(0, max_pages - 1)]

    for palier in paliers:
        url_page = f"{url_listing}?start={palier}&show=&order="
        resp = requests.get(url_page, headers=HEADERS, timeout=15)
        nouvelles = [...]  # extraction + filtre langue, comme ci-dessus
        if not nouvelles:
            break  # fin du listing atteinte
        urls_trouvees += nouvelles
    return urls_trouvees
\end{lstlisting}

Le d\'edoublonnage entre deux ex\'ecutions successives ne se fait pas ici, mais repose sur la contrainte \texttt{document.url UNIQUE} d\'ej\`a en place en base (voir section~7)~: relancer la d\'ecouverte ne recr\'ee jamais un document d\'ej\`a connu.

\subsection{Validation en conditions r\'eelles}
La d\'ecouverte a \'et\'e ex\'ecut\'ee pour de vrai (\texttt{scripts/decouverte\_publications.py}) sur une page listing r\'eelle~: 10/10 URLs d'articles retrouv\'ees, identiques \`a une v\'erification manuelle ind\'ependante de la m\^eme page. Un probl\`eme r\'eel est apparu \`a cette occasion~: une page en version arabe (\texttt{...-version-Ar\_a4217.html}) s'est gliss\'ee dans les r\'esultats. Corrig\'e en r\'eutilisant l'heuristique de d\'etection de langue d\'ej\`a valid\'ee en Sprint~1 pour les pi\`eces jointes (\texttt{\_filtrer\_urls\_arabe}), test\'e avec un cas reconstruit \`a partir de cet exemple r\'eel.

L'inventaire des pages listing a \'et\'e compl\'et\'e en parcourant r\'eellement le site (\texttt{data/listing\_urls.py})~: Population~\&~d\'emographie 10/10 sous-th\`emes, March\'e du travail 1 page agr\'eg\'ee (cas particulier confirm\'e), \'Economie partiellement (4 pages couvrant Comptes nationaux, Conjoncture, Conjoncture entreprise, \'Etudes \'economiques -- reste \`a v\'erifier Indices des prix, Secteurs d'activit\'e, Sph\`ere informelle).

\section{Support des fichiers DOCX (ADR 0005)}
D\'ecouvert en Sprint~1 en testant l'extraction sur de vraies publications~: certaines notes hcp.ma (indices IPC/IPPI notamment) sont diffus\'ees en \texttt{.docx} plut\^ot qu'en PDF/XLSX. D\'ecision prise avec Ayman le 17~juillet~: \'etendre \texttt{Extracteur} plut\^ot que d'ignorer ce format, via \texttt{python-docx}. La d\'etection du type de fichier r\'eel (\texttt{\_contenu\_semble\_valide}) a d\^u \^etre renforc\'ee car DOCX et XLSX sont tous deux des archives ZIP au niveau binaire -- distinction faite sur la pr\'esence de fichiers internes propres \`a chaque format (\texttt{word/document.xml} vs \texttt{xl/workbook.xml}).

\section{Chunking et embeddings}

\subsection{Chunking}
Aucune taille de chunk ni chevauchement n'\'etait fix\'e dans les documents de conception (v\'erifi\'e en relisant \texttt{fiche\_cadrage\_v4.pdf} \S8.1 et \texttt{conception\_uml\_v3.pdf} \S8, qui demandent un d\'ecoupage ``respectant la structure des documents'' sans plus de pr\'ecision). D\'ecision d'impl\'ementation prise et justifi\'ee dans le code~: d\'ecoupage par paragraphe avec empaquetage glouton, taille cible 1500 caract\`eres (~300-400 tokens en fran\c{c}ais), chevauchement 200 caract\`eres pour ne pas couper le contexte entre deux chunks cons\'ecutifs. Approximation par caract\`eres plut\^ot que par le vrai tokenizer BGE-M3, pour ne pas devoir charger le mod\`ele (~2~Go) uniquement pour compter des tokens.

\begin{lstlisting}[caption={\texttt{src/indexeur\_texte.py} -- decoupage par paragraphe, empaquetage glouton}]
TAILLE_CHUNK_CARACTERES = 1500
CHEVAUCHEMENT_CARACTERES = 200

@staticmethod
def _decouper_en_chunks(texte, taille_max=TAILLE_CHUNK_CARACTERES,
                         chevauchement=CHEVAUCHEMENT_CARACTERES):
    paragraphes = [p.strip() for p in texte.split("\n") if p.strip()]
    chunks, courant = [], ""
    for paragraphe in paragraphes:
        if len(paragraphe) > taille_max:
            if courant:
                chunks.append(courant); courant = ""
            pas = max(taille_max - chevauchement, 1)
            for i in range(0, len(paragraphe), pas):
                chunks.append(paragraphe[i : i + taille_max])
            continue
        candidat = f"{courant}\n{paragraphe}" if courant else paragraphe
        if len(candidat) <= taille_max:
            courant = candidat
        else:
            chunks.append(courant)
            amorce = courant[-chevauchement:] if chevauchement else ""
            courant = f"{amorce}\n{paragraphe}" if amorce else paragraphe
    if courant:
        chunks.append(courant)
    return chunks
\end{lstlisting}

Un paragraphe seul plus grand que \texttt{taille\_max} (ex. tableau aplati en un bloc de texte) est tronqu\'e \`a la dure plut\^ot qu'ignor\'e, pour ne perdre aucun contenu.

\subsection{Embeddings (BGE-M3)}
Le mod\`ele retenu par l'ADR 0002 est \texttt{BAAI/bge-m3} (multilingue, via \texttt{sentence-transformers}), confirm\'e comme choix de production dans les documents de conception. La fonction d'embedding est injectable au constructeur d'\texttt{IndexeurTexte}, pour permettre les tests sans t\'el\'echarger le mod\`ele (~2~Go) dans l'environnement de d\'eveloppement sans acc\`es r\'eseau. \textbf{Point de vigilance important}~: cette fonction factice n'est utilis\'ee QUE dans les tests automatis\'es -- le chemin de production (ci-dessous) charge et ex\'ecute le vrai mod\`ele, valid\'e r\'eellement en section~8.

\begin{lstlisting}[caption={\texttt{src/indexeur\_texte.py} -- chargement paresseux du vrai modele}]
def _embarquer(self, textes: list[str]) -> list[list[float]]:
    if self._fonction_embedding is not None:
        return self._fonction_embedding(textes)  # utilise UNIQUEMENT dans les tests
    if self._modele is None:
        from sentence_transformers import SentenceTransformer
        self._modele = SentenceTransformer(NOM_MODELE_EMBEDDING)  # "BAAI/bge-m3"
    return self._modele.encode(textes, normalize_embeddings=True).tolist()
\end{lstlisting}

\section{Indexation hybride~: Chroma (s\'emantique) + BM25 (lexical)}
Chaque chunk est index\'e \`a la fois dans ChromaDB (recherche par similarit\'e d'embeddings, persist\'ee sur disque) et dans un index BM25 (recherche lexicale par mots-cl\'es, via \texttt{rank\_bm25}, reconstruit \`a la demande faute de persistance native de cette biblioth\`eque). Les deux classements sont fusionn\'es par \emph{Reciprocal Rank Fusion} (RRF, $k=60$, valeur usuelle de la litt\'erature -- Cormack, Clarke \& Buettcher, 2009) plut\^ot qu'une moyenne pond\'er\'ee des scores~: les scores denses (cosinus, born\'es entre 0 et 1) et BM25 (non born\'es) ne sont pas sur la m\^eme \'echelle, alors que RRF ne s'appuie que sur le \emph{rang} de chaque r\'esultat dans chaque liste.

\begin{lstlisting}[caption={\texttt{src/indexeur\_texte.py} -- fusion Reciprocal Rank Fusion}]
@staticmethod
def _fusionner_rrf(listes_classees, top_k, k_rrf=60):
    scores = {}
    for liste in listes_classees:
        for rang, id_ in enumerate(liste):
            scores[id_] = scores.get(id_, 0.0) + 1.0 / (k_rrf + rang + 1)
    classement = sorted(scores.keys(), key=lambda i: scores[i], reverse=True)
    return classement[:top_k]
\end{lstlisting}

Point technique v\'erifi\'e exp\'erimentalement pendant le d\'eveloppement~: \texttt{collection.get(ids=[...])} de ChromaDB ne garantit PAS l'ordre des ids demand\'es en retour -- la m\'ethode \texttt{rechercher} r\'e-ordonne donc explicitement les r\'esultats selon le rang RRF plut\^ot que de faire confiance \`a l'ordre renvoy\'e par Chroma.

\section{Persistance SQLite}
Le module \texttt{src/base\_donnees.py} relie les objets Python (\texttt{Document}, \texttt{Chunk}, \texttt{Indicateur}) aux tables de \texttt{db/schema.sql}, avec deux r\`egles de d\'edoublonnage~:

\begin{longtable}{@{}p{4.5cm}p{10.5cm}@{}}
\rowcolor{navy}
\tblhead{Table} & \tblhead{R\`egle de d\'edoublonnage} \\
\midrule
\endhead
\texttt{document} & Contrainte \texttt{url UNIQUE} en base~: \texttt{inserer\_document} renvoie l'id existant si l'URL est d\'ej\`a connue, ne duplique jamais. \\
\texttt{indicateur} & Pas de contrainte UNIQUE dans le sch\'ema~; upsert applicatif sur (nom, p\'eriode, r\'egion, code\_bds) c\^ot\'e code, n\'ecessaire car le pr\'e-remplissage BDS (ADR 0004) est con\c{c}u pour tourner \`a r\'ep\'etition (t\^ache nocturne) sans accumuler de doublons \`a chaque ex\'ecution. \\
\end{longtable}

L'embedding (liste de flottants) est s\'erialis\'e en BLOB via le module standard \texttt{array} (float32), sans d\'ependance suppl\'ementaire~:

\begin{lstlisting}[caption={\texttt{src/base\_donnees.py} -- serialisation de l'embedding}]
def _embedding_vers_blob(embedding: list[float]) -> bytes:
    return array.array("f", embedding).tobytes()

def blob_vers_embedding(blob: bytes) -> list[float]:
    tampon = array.array("f")
    tampon.frombytes(blob)
    return list(tampon)
\end{lstlisting}

\section{Scripts d'orchestration}
\begin{longtable}{@{}p{4.5cm}p{10.5cm}@{}}
\rowcolor{navy}
\tblhead{Script} & \tblhead{R\^ole} \\
\midrule
\endhead
\texttt{scripts/indexer\_documents.py} & Cha\^ine compl\`ete~: d\'ecouverte (\texttt{data/listing\_urls.py}) \`a Extracteur \`a insertion SQLite \`a \texttt{IndexeurTexte.indexer}. Les documents HTML sont explicitement ignor\'es (ADR 0003). Option \texttt{--limite N} ajout\'ee pour permettre un test rapide sur quelques documents avant un run complet. \\
\texttt{scripts/preremplir\_indicateurs\_bds.py} & Pr\'e-remplit les indicateurs cur\'es depuis l'API BDS (ADR 0004)~: fait maintenant le vrai upsert en base (avant le Sprint~2, le script affichait seulement un r\'esum\'e sans \'ecrire en SQLite). \\
\texttt{scripts/inspecter\_index.py} & Outil de v\'erification/d\'emonstration~: affiche un r\'esum\'e lisible de la base (documents, chunks, indicateurs) et, avec \texttt{--recherche "question"}, ex\'ecute une vraie recherche hybride -- pens\'e pour montrer un r\'esultat concret \`a l'encadrante. \\
\end{longtable}

\section{Tests automatis\'es}
Philosophie suivie tout au long du projet~: utiliser de vraies biblioth\`eques dans les tests (\texttt{chromadb}, \texttt{rank\_bm25}, \texttt{sqlite3}, \texttt{python-docx}) plut\^ot que des simulacres, en ne substituant que la partie r\'eellement co\^uteuse \`a ex\'ecuter en local (le mod\`ele d'embedding, ~2~Go, sans acc\`es r\'eseau dans l'environnement de d\'eveloppement).

\begin{longtable}{@{}p{6.5cm}p{2cm}p{6.5cm}@{}}
\rowcolor{navy}
\tblhead{Fichier de tests} & \tblhead{Nombre} & \tblhead{Couvre} \\
\midrule
\endhead
\texttt{test\_indexeur\_texte.py} & 15 & Chunking (limites exactes, chevauchement, tronçonnage), fusion RRF, recherche hybride bout en bout (Chroma + BM25 r\'eels) \\
\texttt{test\_base\_donnees.py} & 9 & Cr\'eation des tables, d\'edoublonnage document/indicateur, round-trip de l'embedding \\
\texttt{test\_indexer\_documents.py} & 4 & Cha\^ine compl\`ete sur un vrai fichier \texttt{.docx} g\'en\'er\'e \`a la vol\'ee, filtrage HTML, non-duplication \\
\rowcolor{lightgrey} \textbf{Total du projet (tous sprints)} & \textbf{52} & \textbf{52/52 verts au 20 juillet 2026} \\
\end{longtable}

\section{Validation en conditions r\'eelles}
Ce sandbox de d\'eveloppement n'a pas d'acc\`es r\'eseau vers hcp.ma ni la capacit\'e de t\'el\'echarger le mod\`ele BGE-M3 (~2~Go)~: toute validation r\'eelle du pipeline complet a donc \'et\'e effectu\'ee par Ayman, sur sa propre machine, le 20~juillet~2026. Les blocs ci-dessous reconstituent fid\`element (texte identique) la sortie r\'eelle de ces ex\'ecutions.

\subsection{Premier test, 3 documents}
\begin{terminal}[title={PS> python -m scripts.indexer\_documents --limite 3}]
\begin{Verbatim}
[Scraper] piece jointe arabe ignoree (hors scope V1) : hcp.ma/file/248485/
[Scraper] piece jointe arabe ignoree (hors scope V1) : hcp.ma/file/245585/
html  | ignore_html  |   0 chunks | Les comptes regionaux. PIB et depenses de...
Warning: You are sending unauthenticated requests to the HF Hub.
(telechargement de BAAI/bge-m3 : modules.json, config.json, pytorch_model.bin
 1.52GB, model.safetensors 1.52GB, tokenizer, sentencepiece.bpe.model ... )
pdf   | ok           |  40 chunks | Les comptes regionaux. PIB et depenses de...
html  | ignore_html  |   0 chunks | Les comptes nationaux provisoires 2025...
Total : 3 document(s) traite(s), 40 chunk(s) indexe(s).
\end{Verbatim}
\end{terminal}

Premi\`ere ex\'ecution r\'eelle de la cha\^ine compl\`ete~: t\'el\'echargement effectif du mod\`ele BGE-M3, scraping r\'eel de hcp.ma, extraction, chunking, indexation. Les 2 pi\`eces jointes arabes sont correctement \'ecart\'ees (filtre de langue), les 2 pages HTML sont correctement ignor\'ees (ADR 0003), et 1 vrai document PDF est extrait, d\'ecoup\'e en 40 chunks et index\'e.

\subsection{V\'erification de la recherche hybride}
\begin{terminal}[title={PS> python -m scripts.inspecter\_index --recherche "produit interieur brut"}]
\begin{Verbatim}
=== Resume de la base SQLite ===
Documents  : 1     Chunks : 40     Indicateurs: 0

=== Extraits de chunks (3 premiers) ===
  chunk #1 (position 0, embedding de dimension 1024) :
    "Table des matieres Contexte economique national 04..."

=== Recherche de demonstration : "produit interieur brut" ===
  1. (document #1, position 29)
     "Produit interieur brut de la region 2,1 2,0 33125 29921 10,7..."
  2. (document #1, position 27)
     "Produit interieur brut (aux prix courants) par branche d'activite..."
  3. (document #1, position 0)
     "Table des matieres Contexte economique national..."
\end{Verbatim}
\end{terminal}

Point de v\'erification cl\'e~: la dimension d'embedding affich\'ee est \textbf{1024}, la signature exacte de BGE-M3 -- preuve que le vrai mod\`ele a tourn\'e, pas la fonction factice utilis\'ee dans les tests. Les 5 r\'esultats retourn\'es (3 montr\'es ici) sont tous effectivement pertinents sur le sujet demand\'e, bien class\'es.

\subsection{Run \'elargi, 15 documents}
\begin{terminal}[title={PS> python -m scripts.indexer\_documents --limite 15}]
\begin{Verbatim}
pdf   | ok |  40 chunks | Les comptes regionaux...
pdf   | ok | 119 chunks | Les comptes nationaux provisoires 2025...
pdf   | ok |  54 chunks | Les comptes nationaux des secteurs institutionnels...
pdf   | ok |  73 chunks | Les provinces du Sud du Royaume...
[Scraper] piece jointe arabe ignoree (hors scope V1) : (x4)
pdf   | ok |  49 chunks | Budget economique exploratoire 2027...
pdf   | ok |  33 chunks | Point de conjoncture N. 50, Juillet 2026
pdf   | ok |  87 chunks | Note de conjoncture N. 48, Avril 2026
Total : 15 document(s) traite(s), 455 chunk(s) indexe(s).
\end{Verbatim}
\end{terminal}

Le mod\`ele BGE-M3, t\'el\'echarg\'e et mis en cache lors du premier test, se charge ici quasi instantan\'ement (pas de nouveau t\'el\'echargement). 8 documents PDF r\'eels indexes sans erreur, 455 chunks au total, filtre de langue toujours actif (4 pi\`eces jointes arabes suppl\'ementaires \'ecart\'ees). Ce lot n'a couvert que la cat\'egorie \'Economie (limite atteinte avant les listings Marché du travail/Population)~: la validation r\'eelle du support DOCX (section~4) reste \`a faire sur un run plus large ou cibl\'e, prochaine \'etape imm\'ediate (voir section~9).

\section{Limites connues et travail restant}
\begin{longtable}{@{}p{6cm}p{9cm}@{}}
\rowcolor{navy}
\tblhead{Point} & \tblhead{\'Etat} \\
\midrule
\endhead
Validation DOCX en conditions r\'eelles & Le support DOCX (ADR 0005) est impl\'ement\'e et test\'e (fixture g\'en\'er\'ee \`a la vol\'ee), mais pas encore rencontr\'e lors d'une ex\'ecution r\'eelle sur hcp.ma -- les runs effectu\'es \`a ce jour n'ont produit que des PDF. \\
Inventaire des pages listing \'Economie & Partiel~: Indices des prix et production (IPC/IPPI/IPI/ICE), Secteurs d'activit\'e et Sph\`ere informelle restent \`a v\'erifier et ajouter \`a \texttt{data/listing\_urls.py}. \\
\texttt{ConstructeurIndicateurs.structurer} & Non impl\'ement\'e (repli PDF/XLSX pour les indicateurs hors catalogue BDS)~: priorit\'e volontairement basse, l'API BDS (ADR 0004) couvre d\'ej\`a la majorit\'e des indicateurs cibl\'es. \\
\rowcolor{lightgrey} Piste \texttt{hcp.ma/downloads/?tag=X} & D\'ecouverte alternative trouv\'ee (base de t\'el\'echargements filtrable par tag, liens directs) mais volontairement non exploit\'ee -- changement d'architecture plus lourd, \`a discuter dans un \'eventuel ADR ult\'erieur. \\
\end{longtable}

\section{Transition vers le Sprint 3}
Le Sprint 2 est fonctionnellement complet c\^ot\'e code, et sa premi\`ere validation en conditions r\'eelles (section~8) confirme que la cha\^ine compl\`ete fonctionne effectivement de bout en bout, avec le vrai mod\`ele BGE-M3 et de vraies publications hcp.ma -- pas seulement en tests avec des composants factices. Reste \`a \'elargir cette validation (couverture DOCX, plus grand volume de documents) avant d'attaquer le Sprint 3~: routage des questions, recherche structur\'ee sur les indicateurs BDS, reranking, et g\'en\'eration de r\'eponse.

\end{document}
